\documentclass[11pt]{article}
% Shared preamble for the research-note series.
% Create note-specific content in the per-note directory; keep common macros here.
\usepackage[margin=1in]{geometry}
\usepackage[T1]{fontenc}
\usepackage[utf8]{inputenc}
\usepackage{amsmath,amssymb}
\usepackage{booktabs,longtable,array}
\usepackage{graphicx}
\usepackage{float}
\usepackage{enumitem}
\usepackage{xcolor}
\usepackage{hyperref}
\usepackage{caption}
\usepackage{titlesec}
\hypersetup{colorlinks=true, linkcolor=blue!60!black, urlcolor=blue!60!black, citecolor=blue!60!black}
\setlist[itemize]{leftmargin=1.25em, itemsep=0.2em, topsep=0.3em}
\setlist[enumerate]{leftmargin=1.4em, itemsep=0.2em, topsep=0.3em}
\setlength{\parindent}{0pt}
\setlength{\parskip}{0.6em}
\newcommand{\statusbox}[2]{%
  \noindent\fcolorbox{black}{gray!10}{%
    \parbox{\dimexpr\linewidth-2\fboxsep-2\fboxrule\relax}{%
      \textbf{Status:} #1\\#2
    }%
  }
}

% Shared notation helpers for the research-note series.
\newcommand{\Jlin}{J}
\newcommand{\Jphys}{J_{\mathrm{phys}}}
\newcommand{\Jsurf}{J_{\mathrm{surf}}}
\newcommand{\JdB}{J_{\mathrm{dB}}}
\newcommand{\phiopt}{\phi_{\mathrm{opt}}}
\newcommand{\sigmathreedb}{\sigma_{3\mathrm{dB}}}
\newcommand{\Nt}{N_t}
\newcommand{\Nphi}{N_{\phi}}
\newcommand{\Lfiber}{L}
\newcommand{\Pcont}{P}
\newcommand{\band}{\mathcal{B}_{\mathrm{Raman}}}
\newcommand{\todoitem}[1]{\textcolor{red!70!black}{\textbf{TODO:} #1}}


\title{NOTE TITLE}
\author{Fiber Raman Suppression Project}
\date{Evidence snapshot: YYYY-MM-DD}

\begin{document}
\maketitle

\begin{abstract}
State the research question, the main conclusion, and the evidence status in
4--6 sentences.  Keep this readable by an advanced undergraduate.  Do not use
internal milestone labels.
\end{abstract}

\section*{Claim Status}
\statusbox{established / provisional / exploratory}{One or two sentences
explaining exactly what is supported and what is not yet supported.}

\section{Question and Thesis}

Define the narrow question this note answers.  Then state the thesis in plain
language.  If there are multiple claims, separate the mature claim from the
provisional claim.

\section{Setup and Common Notation}

Give only the setup needed for this lane: fiber, pulse, control variable,
Raman band, grid family, and output metric.  Reuse notation from the rest of
the note series.

\section{Objective and Optimizer Coordinates}

Write the full scalar objective:
\[
J(\theta)=
\frac{\sum_{k\in\mathcal B} w_k |\hat u_k(L;\theta)|^2}
{\sum_k w_k |\hat u_k(L;\theta)|^2}.
\]
If the plotted objective is dB-scaled, write
\[
J_{\mathrm{dB}}=10\log_{10}J.
\]
Then define the optimizer coordinate map, for example
\[
\hat u_k(0;\phi)=\hat u_{0,k}\exp(i\phi_k),
\qquad
\phi = Bc,
\qquad
a_k=A_k\exp(i\phi_k).
\]
Include regularizers, projections, bounds, clipping, or transformed variables
if the code uses them.

\paragraph{TL;DR.}
Add a short public-facing explanation of what the objective and coordinates
mean.  This should help a reader explain the section without copying private
agent notes.

\section{Math Delta}

Include only the equations unique to this research direction: reduced-basis
chain rules, sharpness penalties, trust-region model, multimode mode sums,
amplitude chain rules, continuation maps, or transfer metrics.

\subsection*{Intuition Check}

Explain the math in words.  If no additional math is needed, say that this
lane uses the common objective and skip this section in the real note.

\section{Implementation Surface}

\begin{center}
\small
\begin{tabular}{p{0.34\linewidth}p{0.56\linewidth}}
\toprule
\textbf{Component} & \textbf{Role} \\
\midrule
\texttt{script\_or\_module.jl} & What this file contributes to the result. \\
\texttt{result\_summary.md} & Where the reported numbers come from. \\
\texttt{standard\_images/} & Figure source for the note. \\
\bottomrule
\end{tabular}
\end{center}

\section{Experimental Strategy}

Describe the comparison set: control, baseline, candidate method, transfer
test, robustness probe, sweep, continuation ladder, or validation gate.  Say
what was excluded and why.

\section{Representative Results}

Start with a compact result table.  Include units, convergence status, and
whether each row is a control, native optimization, transfer evaluation, or
rerun.

\clearpage
\section{No-Optimization Control}

Include the no-shaping/no-optimization control page when the artifact exists.
If there is no control artifact yet, leave a clear missing-evidence sentence
instead of using a placeholder.

\begin{figure}[H]
\centering
% \includegraphics[width=0.48\linewidth]{figures/control_phase_diagnostic.png}
% \includegraphics[width=0.48\linewidth]{figures/control_evolution.png}
\caption{Control/reference case. Explain what this baseline establishes.}
\end{figure}

\clearpage
\section{Representative Figure Pair}

Pair the phase diagnostic/profile with the corresponding evolution heat map on
the same page.

\begin{figure}[H]
\centering
\begin{minipage}{0.48\linewidth}
\centering
% \includegraphics[width=\linewidth]{figures/main_phase_diagnostic.png}
\captionof{figure}{Phase diagnostic/profile for the representative case.}
\end{minipage}\hfill
\begin{minipage}{0.48\linewidth}
\centering
% \includegraphics[width=\linewidth]{figures/main_evolution.png}
\captionof{figure}{Corresponding evolution heat map for the same case.}
\end{minipage}
\end{figure}

\section{Result Provenance}

List the local artifact sources for the reported numbers and figures.  Do not
make readers reverse-engineer which run produced a table row.

\section{Interpretation}

Explain what the result means physically and numerically.  Separate mechanism,
performance, robustness, and compute-cost claims.

\section{Presentation Capsule}

\begin{itemize}
\item \textbf{Slide takeaway:} One sentence the audience should remember.
\item \textbf{Use this figure:} Name the best figure pair or diagram for a
  presentation slide.
\item \textbf{Say out loud:} Two to four short points that explain the method,
  result, limitation, and why it matters.
\item \textbf{Do not overclaim:} One sentence naming the boundary of the
  evidence.
\end{itemize}

\section{Limitations and Missing Evidence}

List unresolved gaps directly: missing reruns, non-convergence, absent control
artifacts, incomplete standard images, unverified gradients, or unsupported
generalizations.

\section{Reproduction Capsule}

Give the command family, machine boundary, key parameters, expected outputs,
and visual-inspection requirement.  If the exact command is still unstable,
say what entry point must be stabilized before the note can become final.

\section{Sources}

Use external research sources for physics and numerical methods.  Use internal
paths only for result provenance.

\end{document}
